
\subsection{The Stability}

We must make sure that the approximation solution obtained form whether the Lax--Friedrichs type scheme or Godunov type scheme shares similar properties as the entropy solution. Since the schemes may devolve undesired oscillations, we must impose a condition on the parameters $\Delta x$, $\Delta t$ and $\alpha$. So, we choose the parameters such that
\begin{equation}\label{eq:stability LF}
    \Delta t \leq \frac{\Delta x}{\alpha +2} \text{ for Lax--Friedrichs}
\end{equation}
and
\begin{equation}\label{eq:stability Gud}
    \Delta t \leq \Delta x \text{ for Godunov}
\end{equation}
in order to obtain stability. In order to understand the stability conditions deeply, see see \textcolor{red}{Citation is needed here}.

\textbf{STEP 1.} Initialize $\mathcal{U}_{j}^{0}$ by setting:
\begin{align*}
    \mathcal{U}_{j}^{0} = \frac{1}{\Delta x}\int_{x_{j-\sfrac{1}{2}}}^{x_{j+\sfrac{1}{2}}} \mathcal{U}_{0}(x)\,dx, \quad \forall j \in \{0,\dots,K-1\}.
\end{align*}
\textbf{STEP 2.}   for $n = 0,\dots,N-1$, compute:
\begin{align*}
    \mathcal{U}_{j}^{n+1} & = \mathcal{U}_{j}^{n} - \frac{\Delta t}{\Delta x}\bigl(F_{j+\sfrac{1}{2}} - F_{j-\sfrac{1}{2}}\bigr)                                                                                                                                                              \\
                          & =  \mathcal{U}_{j}^{n} - \frac{\Delta t}{\Delta x} \Bigl[\frac{1}{2}\bigl(\mathcal{U}_{j+1}^{n}q_{j+1}^{n} - \mathcal{U}_{j-1}^{n}q_{j-1}^{n}\bigr) - \frac{\alpha}{2}\big(\mathcal{U}_{j+1}^{n} + \mathcal{U}_{j-1}^{n}\big) + \alpha \mathcal{U}_{j}^{n}\Bigr].
\end{align*}

The second step is also supposed to hold for all $j \in \{0,\dots,K-1\}$. However, for some nodes for example $j = 0$ and  $j = K-1$ we obtain:
\[
    \mathcal{U}_{0}^{n+1} = \mathcal{U}_{j}^{n} - \frac{\Delta t}{\Delta x} \Bigl[\frac{1}{2}\big(\mathcal{U}_{1}^{n}q_{1}^{n} - \mathcal{U}_{-1}^{n}q_{-1}^{n}\big) - \frac{\alpha}{2}\big(\mathcal{U}_{1}^{n} + \mathcal{U}_{-1}^{n}\big) + \alpha \mathcal{U}_{0}^{n}\Bigr],
\]
and
\[
    \mathcal{U}_{K-1}^{n+1} = \mathcal{U}_{j}^{n} - \frac{\Delta t}{\Delta x} \Bigl[ \frac{1}{2}\big(\mathcal{U}_{K}^{n}q_{K}^{n} - \mathcal{U}_{K-2}^{n}q_{K-2}^{n}\big) - \frac{\alpha}{2}\big(\mathcal{U}_{K}^{n} + \mathcal{U}_{K-2}^{n}\big) + \alpha \mathcal{U}_{K-1}^{n}\Bigr],
\]
respectively.
We observe that for some nodes, the scheme involves undefined terms such as $\mathcal{U}_{-1}^{n}$ and $\mathcal{U}_{K-1+m}^{n}$ for $ m = 0, \cdots M-1$. Therefore, without additional information the scheme is not valid for all nodes.

The above discussion indicate that specifying boundary conditions are needed, so we take a look at how the boundary conditions are handled.
There are several options for specifying the boundary conditions on the spatial domain $[x_L, x_R]$, but the choice depends on whether we have an initial value problem (IVP) or a boundary value problem (BVP).

\begin{itemize}
    \item \textbf{Initial Value Problem (IVP):} Artificial boundary conditions are required to handle the truncated domain $[x_L, x_R]$. Truncation is necessarily for initial value problem on the entire real line $\mathbb{R}$ as computational domain needs to be bounded.
    \item \textbf{Boundary Value Problem (BVP):} Common choices include  Periodic, Dirichlet, and Neumann boundary conditions.
\end{itemize}

In this thesis, we focus on \textbf{artificial boundary conditions} and \textbf{Periodic boundary conditions}. For a detailed discussion of the remaining types, see \textcolor{red}{Citation needed here}.

\textbf{Artificial boundary conditions} are implemented as
\[
    \mathcal{U}_{-1}^{n}= \mathcal{U}_{0}^n \quad \text{and } \quad \mathcal{U}_{K-1+m}^{n} = \mathcal{U}_{K-1}^n \text{ for } m = 0, \cdots M-1,
\]
and \textbf{Periodic boundary conditions} are implemented as
\[
    \mathcal{U}_{-1}^{n}= \mathcal{U}_{K-1}^n \quad \text{and } \quad \mathcal{U}_{K-1+m}^{n} = \mathcal{U}_{m}^n \text{ for } m = 0, \cdots M-1.
\]
For the Godunov flux
\[
    g(\rho_L,\rho_R,q_L,q_R)=\rho_L(1-q_R),
\]
the finite volume update is
\[
    H(\rho_{j-1}^n,\rho_j^n,\ldots,\rho_{j+m}^n)
    =
    \rho_j^n
    +
    \lambda
    \left[
        \rho_{j-1}^n(1-q_j^n)
        -
        \rho_j^n(1-q_{j+1}^n)
        \right],
\]
where
\[
    q_j^n=\sum_{\ell=0}^{m-1}w_\ell \rho_{j+\ell}^n,
    \qquad
    q_{j+1}^n=\sum_{\ell=0}^{m-1}w_\ell \rho_{j+1+\ell}^n.
\]

The partial derivatives of \(H\) are as follows.

For \(k=-1\),
\[
    \frac{\partial H}{\partial \rho_{j-1}^n}
    =
    \lambda(1-q_j^n).
\]

For \(k=0\),
\[
    \frac{\partial H}{\partial \rho_j^n}
    =
    1
    -
    \lambda(1-q_{j+1}^n)
    -
    \lambda w_0\rho_{j-1}^n.
\]
Equivalently,
